% ============================================================================
% Section I: Introduction
% ============================================================================

\section{Introduction}
\label{sec:introduction}

The transition toward low-carbon energy systems requires planning tools
that co-optimize generation, transmission and storage investments while
capturing the operational detail of variable renewables, storage and
hydro cascades~\cite{lumbreras2016progressive}.  Generation and
Transmission Expansion Planning (GTEP) determines the minimum-cost
combination of investments (CAPEX) and operating decisions (OPEX) over
long horizons under
uncertainty~\cite{hemmati2013comprehensive, mahdavi2019tep_survey}.
In hydro-dominated systems the same questions extend to the medium term
(months to years), where reservoir management under inflow uncertainty
is solved with stochastic dual dynamic programming
(SDDP)~\cite{pereira1991sddp}, and to daily operation, where unit
commitment and network constraints set the dispatch and the marginal
prices.

In practice these three horizons are usually served by different
models, with different data formats, element definitions and
simplifications.  A long-term expansion model, a medium-term
hydrothermal model and a short-term scheduling model can therefore
disagree on the value of water, the opportunity cost of storage or the
feasibility of operating rules, and the policy handed from one horizon
to the next is only approximately consistent.  The commercial tools used
by system operators and utilities (e.g.\ PLEXOS, PSR SDDP, and in Chile
the PLP model maintained by the system operator) are closed or tied to
specific data formats, which limits the auditability of their results.
Open-source tools such as PyPSA~\cite{brown2018pypsa},
GenX~\cite{jenkins2017genx}, IDAES-GTEP~\cite{idaes_gtep2024} and
SDDP.jl~\cite{dowson2020sddp_jl} have broadened access to planning and
stochastic optimization.  None of them, however, combines in
one model a detailed hydro and water-rights representation, a network
with losses, SDDP and capacity expansion.

This paper presents GTOpt, an open-source optimizer for planning and
operating energy systems, written in C++26.  GTOpt descends from the
FESOP (Fabulous Energy System Optimizer)
framework~\cite{buitrago2022fesop}, which extended long-term
hydrothermal scheduling~\cite{pereira2016hydrothermal} with capacity
expansion, variable renewables and modern flexibility mechanisms.  Its
central design decision is to use \emph{one} formulation of the energy
system---electricity, water, fuels and LNG, hydrogen, ammonia and
heat---at every time scale, and to change only the time structure and
the solution method between scales.

The main contributions of this work are:

\begin{enumerate}
  \item A complete LP/MIP formulation of the energy system
    (Section~\ref{sec:formulation}).  It includes:
    \begin{itemize}
      \item DC optimal power flow with Kirchhoff constraints,
        transmission switching, zero-impedance ties, dynamic line
        rating through critical spans, and several line-loss models
        with a guard against loss arbitrage under negative prices;
      \item batteries with charge and discharge modes, and hydro
        cascades with reservoirs, turbines, pumps, seepage and
        head-dependent efficiency;
      \item a generic water-rights layer for irrigation agreements;
      \item typical periods and point-sampled stages on Chebyshev
        nodes, with storage balances consistent with the quadrature of
        energy and cost;
      \item unit commitment (three-binary energy-block and power-based
        formulations on chronological stages, a per-block formulation
        on load-duration stages) with
        piecewise heat rates and pricing modes, dynamic reserve
        requirements, battery reserve deliverability and synchronous
        inertia;
      \item LNG terminals and hydrogen, ammonia and heat chains;
      \item a natural-gas network with Weymouth pipes in squared
        pressure, elevation, compressor fuel, line pack, and
        gas-fired units limited by their fuel-gas pressure, with
        nodal gas prices (Section~\ref{sec:gas});
      \item multi-pollutant emission zones with caps, banking and
        carbon capture;
      \item multi-stage capacity expansion of generation, transmission,
        storage and demand.
    \end{itemize}

  \item A multi-scale design in which the same element definitions are
    used from 15-minute operation to 30-year planning, with the scales
    coupled through SDDP optimality and feasibility cuts
    (Section~\ref{sec:multiscale}).

  \item Three solution methods---monolithic LP/MIP; SDDP with
    apertures, multi-cut and elastic-filter feasibility cuts; and a
    multi-level cascade---plus an investment master over SDDP operation
    (Section~\ref{sec:methods}).

  \item An open software architecture (Section~\ref{sec:architecture}):
    \begin{itemize}
      \item JSON and Apache Parquet data, and PAMPL user constraints;
      \item run-time solver plugins and solver-neutral warm starts;
      \item exportable LP files with structural fingerprints;
      \item converters from PLP, PLEXOS, PSR SDDP, PSS/E, pandapower
        and open test-system formats;
      \item command-line, terminal-dashboard, browser and web-service
        interfaces.
    \end{itemize}

  \item Validation on IEEE benchmark networks, pandapower, published
    SDDP, unit-commitment, pricing and reserve benchmarks, and open gas
    and gas--electric test systems (Section~\ref{sec:validation}), and
    a comparison of
    GTOpt with the official daily schedule of the Chilean national
    system on 14 days, computed by the system operator with PLEXOS
    (Section~\ref{sec:cases}).
\end{enumerate}

The remainder of this paper is organized as follows.
Section~\ref{sec:literature} reviews the state of the art.
Section~\ref{sec:formulation} presents the formulation, and
Section~\ref{sec:gas} its natural-gas network.
Section~\ref{sec:methods} describes the solution methods and the
multi-scale coupling.  Section~\ref{sec:architecture} details the
software architecture.  Section~\ref{sec:validation} presents the
validation on benchmarks, and Section~\ref{sec:cases} the comparison
with the Chilean daily schedule.  Section~\ref{sec:conclusions}
concludes and outlines future work.
